\documentclass[11pt,a4paper]{article}

% ========================
% ESSENTIAL PACKAGES
% ========================

% Language and encoding
\usepackage[utf8]{inputenc}
\usepackage[T1]{fontenc}
\usepackage[english]{babel}
%\usepackage{natbib}

% Math packages
\usepackage{amsmath}        % Enhanced math environments
\usepackage{amssymb}        % Additional math symbols
\usepackage{amsthm}         % Theorem environments
\usepackage{amsfonts}       % Additional math fonts
\usepackage{mathtools}      % Extension to amsmath
\usepackage{bm}             % Bold math symbols
\usepackage{dsfont}         % Double-struck fonts (for sets like ℝ, ℕ)
\usepackage{mathrsfs}       % Script fonts for math

% Graphics and figures
\usepackage{graphicx}       % Include graphics
\usepackage{float}          % Better float placement
\usepackage{subcaption}     % Subfigures
\usepackage{tikz}           % Drawing diagrams
\usepackage{pgfplots}       % Plotting
\pgfplotsset{compat=1.18}

% Page layout and formatting
\usepackage[margin=1in]{geometry}
\usepackage{fancyhdr}       % Custom headers and footers
\usepackage{titlesec}       % Customize section titles
\usepackage{enumitem}       % Customize lists
\usepackage{parskip}        % Paragraph spacing

% Colors and highlighting
\usepackage{xcolor}
\usepackage{mdframed}       % Framed environments
\usepackage{tcolorbox}      % Colored boxes

% Tables
\usepackage{array}          % Extended array environments
\usepackage{booktabs}       % Professional tables
\usepackage{multirow}       % Multi-row cells

% References and links
\usepackage{hyperref}       % Hyperlinks
\usepackage{cleveref}       % Smart cross-references

% Additional useful packages
\usepackage{cancel}         % Cancel terms in equations
\usepackage{siunitx}        % SI units
\usepackage{listings}       % Code listings (if needed)

\usepackage{authblk}

\usepackage[authoryear]{natbib}

% ========================
% CUSTOM COMMANDS
% ========================

% Common math sets
\newcommand{\R}{\mathbb{R}}
\newcommand{\N}{\mathbb{N}}
\newcommand{\Z}{\mathbb{Z}}
\newcommand{\Q}{\mathbb{Q}}
\newcommand{\C}{\mathbb{C}}

% Common operators
\DeclareMathOperator{\lcm}{lcm}
\DeclareMathOperator{\Tr}{Tr}
\DeclareMathOperator{\rank}{rank}
%\DeclareMathOperator{\span}{span}
%\DeclareMathOperator{\null}{null}

% Shortcuts for common expressions
\newcommand{\abs}[1]{\left|#1\right|}
\newcommand{\norm}[1]{\left\|#1\right\|}
\newcommand{\inner}[2]{\left\langle #1, #2 \right\rangle}
\newcommand{\ceil}[1]{\left\lceil #1 \right\rceil}
\newcommand{\floor}[1]{\left\lfloor #1 \right\rfloor}

% Derivatives and integrals
\newcommand{\dd}[1]{\,\mathrm{d}#1}
\newcommand{\dv}[2]{\frac{\mathrm{d}#1}{\mathrm{d}#2}}
\newcommand{\pdv}[2]{\frac{\partial #1}{\partial #2}}

\renewcommand{\hat}{\widehat}
\renewcommand{\tilde}{\widetilde}
\renewcommand{\bar}{\overline}

\newcommand{\vct}{\boldsymbol}
\newcommand{\tr}{\mathrm{tr}}

% ========================
% THEOREM ENVIRONMENTS
% ========================

\theoremstyle{definition}
\newtheorem{definition}{Definition}[section]
\newtheorem{example}{Example}[section]
\newtheorem{exercise}{Exercise}[section]

\theoremstyle{plain}
\newtheorem{theorem}{Theorem}[section]
\newtheorem{lemma}{Lemma}[section]
\newtheorem{corollary}{Corollary}[section]
\newtheorem{proposition}{Proposition}[section]
\newtheorem{assumption}{Assumptio}[section]

\theoremstyle{remark}
\newtheorem{remark}{Remark}[section]
\newtheorem{note}{Note}[section]

\newcommand{\supp}{\mathrm{supp}}
\newcommand{\op}{\mathrm{op}}

\newcommand{\mX}{\mathcal X}
\newcommand{\mD}{\mathcal D}
\newcommand{\bP}{\mathbb P}
\newcommand{\lp}{L^2(\bP)}
\newcommand{\lx}{L^2(\mX)}
\newcommand{\ltwo}{L^2(\mX)}
\newcommand{\ud}{\mathrm d}
\newcommand{\mH}{\mathcal H}
\newcommand{\mA}{\mathcal A}

\newcommand{\diag}{\mathrm{diag}}

\newcommand{\mI}{\mathcal I}

\newcommand{\iso}{\mathsf{iso}}
\newcommand{\ani}{\mathsf{ani}}
\newcommand{\anitr}{\mathsf{anitr}}
\newcommand{\adptr}{\mathsf{adptr}}

\newcommand{\und}{\underline}

% ========================
% COLORED BOXES
% ========================

% Important theorem box
\newtcolorbox{importantbox}{
	colback=blue!5!white,
	colframe=blue!75!black,
	title=Important
}

% Warning/Note box
\newtcolorbox{notebox}{
	colback=yellow!10!white,
	colframe=orange!75!black,
	title=Note
}

% Example box
\newtcolorbox{examplebox}{
	colback=green!5!white,
	colframe=green!75!black,
	title=Example
}

% ========================
% DOCUMENT SETUP
% ========================

% Page style
\pagestyle{fancy}
\fancyhf{}
\rhead{\thepage}
\lhead{\leftmark}


% Hyperlink setup
\hypersetup{
	colorlinks=true,
	linkcolor=blue,
	filecolor=magenta,      
	urlcolor=cyan,
	citecolor=red
}

\title{Honest and Adaptive Pointwise Confidence Intervals in Linear Functional Regression: An Anisotropic Kernel Ridge Regression Approach}
\author[1]{Paromita Dubey \footnote{Author names listed in alphabetical order.}}
\author[2]{Zhiyuan Tang}
\author[2]{Yining Wang}
\affil[1]{Marshall Business School, University of Southern California}
\affil[2]{Naveen Jindal School of Management, University of Texas at Dallas}


\begin{document}

\maketitle

\begin{abstract}
Abstract here.
\end{abstract}

\section{Introduction}

We consider the classical model of linear functional regression with data $\{X_i,Y_i\}_{i=1}^n$, where $X_1,\cdots,X_n\sim P_X$ are independently and identically distributed \emph{functions} 
whose domain is a non-empty compact interval $I\subseteq\mathbb R$, and $\{Y_i\}_{i=1}^n$ are scalar response variables, modeled by
\begin{equation}
\mathbb E[Y_i|X_i] = \eta_0(X_i) := \gamma_0 + \int_I X_i(t)\beta_0(t)\ud t,
\label{eq:model}
\end{equation}
where $\gamma_0\in\mathbb R$ is the intercept term and $\beta_0: I\to\mathbb R$ is the linear functional, both being unknown parameters.
Many well-known methods have been proposed to estimate $\gamma_0$ and $\beta_0$ from data, such as functional principal component analysis (fPCA) approaches \citep[Chapter 8]{ramsay2005functional},
\citep{cai2006prediction} and methods based on reproducing kernel Hilbert spaces \citep{yuan2010reproducing,cai2012minimax}.

In this paper, we consider the question of constructing \emph{pointwise} confidence intervals on $\eta_0(x)$ which are valuable in quantifying the uncertainty in estimates of the functional model
when it is applied to a future data curve $x$. We study the construction of \emph{honest} confidence intervals, which are functions $\varphi$ mapping a data curve $x:I\to\mathbb R$ to a compact interval $\varphi(x)\subseteq\mathbb R$
such that
\begin{equation}
\liminf_{n\to\infty}\inf_{\eta_0\in\Theta} \Pr\left[\eta_0(x)\in\varphi(x)\right] \geq 1-\alpha,\;\;\;\;\;\;\forall x\in\mathrm{supp}(P_X).
\label{eq:honest-ci}
\end{equation}
where $\alpha\in(0,1)$ is a pre-determined significance level,
$\supp(P_X)$ is the support of $P_X$, $\Theta$ is a fixed and known function class and the randomness in Eq.~(\ref{eq:honest-ci}) is over $\{(X_i,Y_i)\}_{i=1}^n$.
To understand how \emph{adative} an honest confidence interval is, we further study the worst-case mean-squared widths of constructed intervals
\begin{equation}
\sup_{\eta_0\in\Theta}\mathbb E_{\eta_0} \left[\int_{\supp(P_X)}\big|\varphi(x)\big|^2\ud P_X(x)\right],
\label{eq:adaptive}
\end{equation}
where $|\varphi(x)|$ is the Lesbesgue measure of the output interval $\varphi(x)$. We are particularly interested in how the above worst-case risk converges to zero as $n\to\infty$,
and how it compares with the mean-squared prediction errors
$$
\sup_{\eta_0\in\Theta}\mathbb E_{\eta_0} \left[\int_{\supp(P_X)}\int_I [x(t)(\hat\eta_n(t)-\eta_0(t)]^2\ud t\ud P_X(x)\right]
$$
for classical fPCA or RKHS estimators $\hat\eta$ of $\eta_0$.

\subsection{Our contributions}

\subsection{Related works}

\subsection{Notations}

For two sequences of real numbers $\{a_n\}$ and $\{b_n\}$, we write $a_n\lesssim b_n$, or equivalently $a_n=O(b_n)$, if there exists a numerical constant $C$ such that $|a_n|\leq C|b_n|$ for all $n$.
We write $a_n\gtrsim b_n$ or $a_n=\Omega(b_n)$ if $b_n\lesssim a_n$, and $a_n\asymp b_n$ or $a_n=\Theta(b_n)$ if both $a_n\lesssim b_n$ and $a_n\gtrsim b_n$ hold.
We denote $L_2(I):=\{x: I\to\mathbb R: \int_I x(t)^2\ud t<+\infty\}$ as the class of all square integrable functions on $I$.
For two functions $x,z\in L_2(I)$, $\langle x,z\rangle$ denotes the standard inner product of $L_2(I)$; that is, $\langle x,z\rangle = \int_I x(t)z(t)\ud t$.

\section{Functional regression via kernel ridge regression}

In this section, we give an overview of an reproducing kernel Hilbert space (RKHS) approach towards linear functional regression, following the seminal works of \cite{cai2012minimax} which also lays the foundation of 
methodological and theoretical contributions of this paper.
To clarify, all results in this section are relatively well-known in the literature, and we are summarizing them for the purpose of completeness only.

Let $K:I\times I\to\mathbb R$ be a real, symmetric, square integrable, and non-negative definite function, known as the \emph{kernel function}.
There then exists a reproducing kernel Hilbert space $\mH\subseteq L_2(I)$ associated with inner product $\langle\cdot,\cdot\rangle_{\mH}$
such that the following reproducing property holds for any $f\in\mH$:
$$
f(t) = \langle K(t,\cdot), f\rangle_{\mH}\;\;\;\;\;\;\forall t\in I.
$$
Assuming $\beta_0\in\mH$. Then a standard kernel ridge regression (KRR) approach to estimate $\alpha_0,\beta_0$ from $\{X_i,Y_i\}_{i=1}^n$ is to solve the following penalized optimization problem
\begin{equation}
\hat\gamma_n,\hat\beta_n \in \arg\min_{\gamma\in\mathbb R,\beta\in\mH} \left\{\frac{1}{n}\sum_{i=1}^n (Y_i-\gamma-\langle X_i,\beta\rangle)^2 + \lambda_n \|\beta\|_\mH^2\right\},
\label{eq:krr}
\end{equation}
where $\lambda_n>0$ is a penalization parameter.
Let $\bar X=\frac{1}{n}\sum_{i=1}^n X_i$ and $\bar Y=\frac{1}{n}\sum_{i=1}^n Y_i$ be the sample averages. It is easy to verify that $\hat\gamma_n=\bar Y$ and 
\begin{equation}
\hat\beta_n\in\arg\min_{\beta\in\mH}\left\{\frac{1}{n}\sum_{i=1}^n (Y_i'-\langle X_i',\beta\rangle)^2 + \lambda_n\|\beta\|_\mH^2\right\}
\label{eq:krr-nointercept}
\end{equation}
where $Y_i'=Y_i-\bar Y$ and $X_i'=X_i-\bar X$.

\subsection{Reduction to an optimization problem in $L_2(I)$}

The work of \cite{cai2012minimax} makes the observation that Eq.~(\ref{eq:krr-nointercept}) can be equivalently reformulated into an infinite-dimensional optimization problem
involving the standard inner product and norm of $L_2(I)$ only, which is helpful for implementation and theoretical analysis purposes.
We summarize here how such equivalence is built below.

By Mercer theorem, there exists an orthonormal basis $\{\phi_j\}_{j=1}^{\infty}$ of $L_2(I)$
and a sequence of non-increasing, strictly positive real numbers $\rho_1\geq\rho_2\geq\cdots$ such that $\phi_j\in\mH$ for all $j$, and furthermore
$$
\langle\phi_j,\phi_k\rangle = \int_I \phi_j(t)\phi_k(t)\ud t = \delta_{jk},\;\;\;\;\;\langle\phi_j,\phi_k\rangle_{\mH}= \frac{\delta_{jk}}{\rho_j},\;\;\;\;\;\;\int_I K(t,\cdot)\phi_j(t)\ud t=\rho_j\phi_j.
$$
where $\delta_{jk}=\vct 1\{j=k\}$ is the Kronecker delta notation.
Note that because $\mH$ is known, $\{\phi_j,\rho_j\}_{j=1}^{\infty}$ are also known.
Following notations in \citep{cai2012minimax}, we define linear operator $L_{K^{1/2}}:L^2(I)\to L^2(I)$ and its inverse operator $L_{K^{-1/2}}$ as follows:
$$
L_{K^{1/2}}(\phi_j) = \sqrt{\rho_j}\phi_j,\;\;\;\;\;\; L_{K^{-1/2}}(\phi_j) = \frac{1}{\sqrt{\rho_j}}\phi_j,\;\;\;\;\;\forall j\geq 1.
$$
In other words, for any $f=\sum_{j=1}^{\infty}\alpha_j\phi_j\in L^2(I)$ such that $\sum_j\alpha_j^2<+\infty$, $L_{K^{1/2}}(f)=\sum_{j=1}^{\infty}\alpha_j\sqrt{\rho_j}\phi_j$
and therefore $L_{K^{1/2}}f\in\mH$ because $\sum_{j=1}^{\infty}\langle L_{K^{1/2}}f,\phi_j\rangle^2/\rho_j <+\infty$.
In fact, \citep[Proposition 2]{yuan2010reproducing} shows that for $f\neq 0$ and $L_{K^{1/2}}f\neq 0$, $f\in L^2(I)$ if and only if $L_{K^{1/2}}f\in\mH$ and furthermore $\int_I f(t)^2\ud t=\|f\|_2^2=\|L_{K^{1/2}}f\|_\mH^2$.
The optimal solution $\hat\beta_n$ in Eq.~(\ref{eq:krr-nointercept}) can then be expressed as $\hat\beta_n = L_{K^{1/2}}\hat\theta_n$, where
\begin{equation}
\hat\theta_n \in \arg\min_{\theta\in L_2(I)}\left\{\frac{1}{n}\sum_{i=1}^n (Y_i'-\langle L_{K^{1/2}}X_i', \theta\rangle^2 + \lambda_n\|\theta\|_2^2\right\},
\label{eq:krr-nointercept-l2}
\end{equation}
where $\|\theta\|_2^2 =\int_I \theta(t)^2\ud t$ is the standard squared $L_2$ norm on $I$.
For brevity, we will also use the notation of $Z_i' = L_{K^{1/2}}X_i'$ throughout this paper.

\subsection{Kernel ridge regression via spectral decomposition}

The infinite-dimensional optimization problem in Eq.~(\ref{eq:krr-nointercept-l2}) can be reduced to a finite-dimensional problem and efficiently solved
using Gram matrices, known as the kernel trick. In this section, we show an alternative characterization of $\hat\theta_n$, representing it as a smoothed version of functional PCA
which is useful motivating our anisotropic estimators later.

Let
$$
T_n := \frac{1}{n}\sum_{i=1}^n Z_i'\otimes Z_i'
$$
be the sample covariance operator, where $Z_i'=L_{K^{1/2}}X_i'$. By the spectral theorem, there exists orthonormal $\{\hat\psi_j\}_{j=1}^n$ in $L_2(I)$ and $\hat s_1\geq \hat s_2\geq\cdots\geq \hat s_n\geq 0$ such that
$$
T_n = \sum_{j=1}^n \hat s_j\hat \psi_j\otimes\hat \psi_j.
$$
The estimator $\hat\theta_n$ can then be expressed as
\begin{equation}
\hat\theta_n = (T_n+\lambda_n I)^{-1}\frac{1}{n}\sum_{i=1}^n Y_i'Z_i' = \frac{1}{n}\sum_{i=1}^n Y_i'\sum_{j=1}^{\infty}\frac{\hat\nu_j(Z_i')}{\hat s_j+\lambda}\hat\psi_j,
\label{eq:krr-nointercept-fpca}
\end{equation}
where $\hat\nu_j(Z_i') = \langle Z_i', \hat\psi_j\rangle=\int_I Z_i'(t)\hat\psi_j(t)\ud t$.

\section{Assumptions}

In this section we list and discuss assumptions that we impose throughout this paper. For some results, additional assumptions are needed: 
we will remark on them when such additional assumptions are introduced specifically for selected theorems or lemmas.

We first introduce mild regularity conditions to make sure the problem is well-posed.
\begin{enumerate}
\item[(A1)] Noise variables $\varepsilon_i:=Y_i-\eta_0(X_i)$ are independently and identically distributed, centered random variables with $\mathbb E[\varepsilon_i^2]=\sigma_\varepsilon^2<+\infty$;
\item[(A2)] $\beta_0\in\mH$ and $\|\beta_0\|_\mH\leq 1$, which is equivalent to $\|\theta_0\|_2=\|L_{K^{-1/2}}\beta_0\|_2 =(\int_I \theta_0(t)^2\ud t)^{1/2} \leq 1$.
\item[(A3)] For $x\sim P_X$, $\|L_{K^{1/2}}x\|_2= (\int_I (L_{K^{1/2}}x)(t))^2\ud t)^{1/2}\leq 1$ almost surely. We also assume that $x\neq 0$ and $L_{K^{1/2}}x\neq 0$ almost surely, to avoid degenerate situations.
\end{enumerate}

We remark that in Assumptions (A2) and (A3) we use one as the upper bounds to $\|\beta_0\|_\mH$ and $\|x\|_2$ for convenience only. Our constructed honest confidence bands do not require such knowledge
and they remain valid with 1 being replaced by any other numerical constants.

We use $\Theta_2(\mH)$ to denote the class of all regression functionals $\beta_0$ that satisfy (A2). Note that $\Theta_2(\mH)$ depends on the selection of the RKHS $\mH$, but not on the covariate distribution $P_X$.

Let $Z:=L_{K^{1/2}}X$ be a random function and $P_Z$ be the distribution of $Z$, which is induced by both $P_X$ and $K$. Consider the population covariance operator
$$
T := \mathbb E_{P_Z}\left[Z\otimes Z\right].
$$
By the spectral theorem, there exists an orthonormal basis $\{\psi_j\}_{j=1}^{\infty}$ of $L_2(I)$ and $s_1\geq s_2\geq\cdots\geq 0$ such that
\begin{equation}
T = \sum_{j=1}^{\infty} s_j\psi_j\otimes\psi_j.
\label{eq:specT}
\end{equation}
Note that, in general, because the population covariance operator $T$ is unknown, $\{\psi_j,s_j\}$ are also unknown.

The following assumptions on imposed directly on the structure of $\{\psi_j\}_j$ and $\{s_j\}_j$.
\begin{enumerate}
\item[(B1)] There exists numerical constants $r>0$ and $0<c_s\leq C_s<+\infty$ such that $c_s j^{-2r}\leq s_j\leq C_sj^{-2r}$ for all $j\geq 1$;
this also implies that $T$ is invertible;
\item[(B2)] There exists a numerical constant $C_p<+\infty$ such that for every $f\in L_2(I)$, $\mathbb E[\langle Z,f\rangle^4] \leq C_p(\mathbb E[\langle Z,f\rangle^2])^2$;
\item[(B3)] There exists numerical constants $\kappa\in[0,1/2)$ and $C_\nu<+\infty$ such that for every $j\geq 1$, $|\langle Z,\psi_j\rangle| \leq C_\nu n^{\kappa}\cdot \sqrt{s_j}$ almost surely.
\end{enumerate}

Assumption (B2) is also imposed, in the form of Assumption (b) in \citep{yuan2010reproducing}. It is satisfied when $X$ is a Gaussian process.

We remark that, in Assumption (B3), the smaller $\kappa$ the stronger the assumption is. For the strongest version of $\kappa=0$, it essentially assumes $|\langle Z,\psi_j\rangle|\lesssim j^{-r}$,
which coincides with coefficient decay assumptions adopted in many fPCA analysis.
{\color{red}In particular, the results in \cite{cai2006prediction} shows that in this case ($\kappa=0$ in (B3)) fPCA achieves an MSE rate of $n^{-(2r-1)/(2r+1)}$ assuming no decay of $|\langle\theta_0,\psi_j\rangle|$.}

To analyze truncated estimators, we also introduce the following assumption which is strictly weaker than Assumption (B3):
\begin{enumerate}
\item[(B3')] For some $q\in(0,1/2)$, there exists numerical constants $\kappa<1/2$ and $C_\nu<+\infty$ such that for any $\delta>0$, there exists $n_0\in\mathbb N$ depending on $\delta$ such that
for every $n\geq n_0$ and $1\leq j\leq n^{q+\delta}$, $|\langle Z,\psi_j\rangle|\leq C_\nu n^{\kappa}\cdot\sqrt{s_j}$.
\end{enumerate}

Given Assumption (B1), the following condition is a sufficient condition that implies Assumption (B3'): for some $a<\kappa/q$, it holds almost surely that
$$
\big|\langle Z,\psi_j\rangle\big|\lesssim j^{-r+a}\;\;\;\;\;\;\forall j\geq 1.
$$

\section{Under-smoothed isotropic kernel regression}\label{sec:isotropic}

As a warm-up exercise and a backup method with minimal assumptions, we first consider honest confidence intervals constructed by \emph{under-smoothing} an isotropic kernel ridge estimator.
Let $\{\lambda_n\}_{n=1}^{\infty}$ be a sequence of positive penalty parameters that satisfies
\begin{equation}
\lambda_n n \to 0,\;\; \lambda_n n^{1-\delta}\to \infty\;\;\;\;\;\;\text{as}\;\;n\to\infty
\label{eq:lambdan-1}
\end{equation}
for any $\delta>0$.
Let $\hat\gamma_n = \bar Y$ and $\hat\beta_n^\iso=L_{K^{1/2}}\hat\theta_n$ where 
$$
\hat\theta_n^\iso = (T_n+\lambda_n I)^{-1}\frac{1}{n}\sum_{i=1}^n Y_i'Z_i',
$$
with $Y_i'=Y_i-\bar Y$, $Z_i'=L_{K^{1/2}}(X_i-\bar X)$ and $T_n=\frac{1}{n}\sum_{i=1}^n Z_i'\otimes Z_i'$.
This is an \emph{under-smoothed} kernel ridge regression estimator with $\lambda_n=o(n^{-1})$, while the optimal penalty parameter choice for estimating $\eta_0$ in mean-squared errors under $P_X$
is $\lambda_n\asymp n^{-2r/(2r+1)}$, per results in \cite{cai2012minimax}.

For a test function $x$ in the support $P_X$, an honest and adaptive confidence interval of $\eta(x)$ at significance level $\alpha\in(0,1)$ is constructed as
$$
\varphi^\iso(x) = [\hat\eta_n^\iso(x)\;\pm\;\tau_{\alpha/2}\hat\sigma_n^\iso(L_{K^{1/2}}x)]
$$
where $\tau_{\alpha/2}$ is the $(1-\alpha/2)$-quantile of the standard Normal distribution, and
$$
\hat\eta_n^\iso(x) = \hat\gamma_n + \int_I x(t)\hat\beta_n^\iso(t)\ud t,\;\;\;\;\;\;\hat\sigma_n(z)^\iso = \frac{\sigma_\varepsilon}{\sqrt{n}}\sqrt{\langle z, (T_n+\lambda_n I)^{-1}z\rangle}.
$$
%To facilitate our analysis, we also introduce the following quantity $\tilde z_n(z)$ defined as
%$$
%\tilde\sigma_n(z) := \frac{\sigma_\varepsilon}{\sqrt{n}}\sqrt{\langle z, (T_n+\lambda_n I)^{-1}T_n(T_n+\lambda_n I)^{-1}z\rangle},
%$$
%which is smaller than $\hat\sigma_n(z)$ almost surely.

\subsection{Honesty of confidence intervals}

For the purpose of our analysis, we also define
we also introduce the following quantity $\tilde z_n(z)$ defined as
$$
\tilde\sigma_n(z) := \frac{\sigma_\varepsilon}{\sqrt{n}}\sqrt{\langle z, (T_n+\lambda_n I)^{-1}T_n(T_n+\lambda_n I)^{-1}z\rangle},
$$
which is smaller than $\hat\sigma_n^\iso(z)$ almost surely.

\begin{theorem}
Suppose Assumptions (A1), (A2), (A3), (B1), (B2), (B3) hold. Suppose also that Eq.~(\ref{eq:lambdan-1}) holds. Then 
for any $x\in\supp(P_X)$, $\hat\eta_n^\iso(x)-\eta_0(x)=b_n(x)+v_n(x)$ such that, as $n\to\infty$, 
\begin{align*}
\frac{b_n(x)}{\hat\sigma_n^\iso(z)}\;\;\overset{p}{\to}\;\; 0,\;\;\;\;\;\;\frac{v_n(x)}{\tilde\sigma_n(z)}\;\;\overset{d}{\to}\;\;\mathcal N(0,1),
\end{align*}
where $z=L_{K^{1/2}}x$.
Furthermore, if $\{\varepsilon_i\}_{i=1}^n$ are Normally distributed, then the above inequalities hold without assuming (B3), and they hold \emph{uniformly} over all $x\in\supp(P_X)$.
\label{thm:isokrr-limiting}
\end{theorem}

The following corollary is then immediate, noting that $\tilde\sigma_n(z)\leq\hat\sigma_n(z)$ almost surely.
\begin{corollary}
Impose conditions in Theorem \ref{thm:isokrr-limiting}. Then
$$
\liminf_{n\to\infty}\inf_{\beta_0\in\Theta_2(\mH)}\Pr[\eta_0(x)\in\varphi^\iso(x)]\;\;\geq\;\; 1-\alpha,\;\;\;\;\;\;\forall x\in\supp(P_X).
$$
Furthermore, if $\{\varepsilon_i\}_{i=1}^n$ are Normally distributed, then the following strengthened result holds and it holds without assuming (B3):
$$
\liminf_{n\to\infty}\inf_{x\in\supp(P_X)}\inf_{\beta_0\in\Theta_2(\mH)}\Pr[\eta_0(x)\in\varphi^\iso(x)] \;\;\geq\;\; 1-\alpha.
$$
\label{cor:isokrr-ci}
\end{corollary}

\subsection{Rates of convergence}

\begin{theorem}
Impose conditions in Theorem \ref{thm:isokrr-limiting}, including Assumption (B3) with some $\kappa<1/2$. Then as $n\to\infty$, 
\begin{equation}
\sup_{\beta_0\in\Theta_2(\mH)}\mathbb E_{\eta_0}\left[\int_{\supp(P_X)}\big|\varphi^\iso(x)\big|^2\ud P_X(x)\right]\lesssim n^{-\frac{2r-1}{2r}+\delta}
\label{eq:simple-validity}
\end{equation}
for any $\delta>0$.
\label{thm:isokrr-rates}
\end{theorem}


\subsection{Minimax optimality}

\begin{theorem}
Fix arbitrary $\alpha\in(0,1)$, $r>0$ and $\kappa=0$. There exists a distribution $P_X$ satisfying Assumptions (A1), (A3), (B1), (B2), and (B3), such that for any $\varphi$ satisfying Eq.~(\ref{eq:honest-ci})
with respect to the constructed $P_X$ and $\Theta_2(\mH)$, it holds that
$$
\sup_{\beta_0\in\Theta_2(\mH)} \mathbb E_{\beta_0}\left[\int_{\supp(P_X)}\big|\varphi(x)\big|^2\ud P_X(x)\right] \gtrsim n^{-\frac{2r-1}{2r}}.
$$
\label{thm:isokrr-lb}
\end{theorem}

\section{Under-smoothed anisotropic kernel regression}

We consider anisotropic extensions to the classical (isotropic) kernel ridge regression, by considering the following estimator
$$
\hat\theta_n^\ani = (T_n+\Lambda_n)^{\dagger}\frac{1}{n}\sum_{i=1}^n Y_i' Z_i'
$$
where $\Lambda_n$ is a positive-semidefinite, self-adjoint covariance operator computed on data that is not necessarily (multiples) of the identity operator,
and $C^\dagger$ is the Moore-Penrose pseudo-inverse of a not necessarily invertible linear operator $C$ (see, e.g.~\citep{fongi2023moore} and references therein).
Because $\Lambda_n$ must be entirely calculated from data, we consider the following form of $\Lambda_n$, rendering it a finite-rank operator:
$$
\Lambda_n = \sum_{j=1}^n \lambda_{nj}\hat\psi_j\otimes\hat\psi_j
$$
where $\{\lambda_{nj}\}_{j=1}^n\geq 0$ only on $\{Z_i'\}_{i=1}^n$ but not $Y_i'$.
With this choice, $\hat\theta_n^\ani$ can be written as
$$
\hat\theta_n^\ani = \frac{1}{n}\sum_{i=1}^n\sum_{j=1}^n \frac{Y_i'\langle Z_i',\hat\psi_j\rangle}{\hat s_j+\lambda_{nj}}\hat\psi_j.
$$

For the purpose of this section, we adopt the choice of 
$$
\lambda_{nj} = \sqrt{\hat s_j\lambda_n}.
$$
With the estimator $\hat\eta^\ani_n(\cdot)$ such that $\hat\eta^\ani_n(x)=\hat\gamma_n+\int_I (L_{K^{1/2}}x)(t)\hat\eta^\ani_n(t)\ud t$, confidence intervals are constructed as
$$
\varphi^\ani(x) := [\hat\eta^\ani_n(x)\;\pm\;\tau_{\alpha/2}\hat\sigma_n^\ani(L_{K^{1/2}}x)]
$$
where $\tau_{\alpha/2}$ is the $(1-\alpha/2)$ quantile of the standard Normal distribution, and $\hat\sigma_n^\ani(\cdot)$ 
$$
\hat\sigma_n^\ani(z) = \frac{\sigma_\varepsilon}{\sqrt{n}}\sqrt{\langle z, (T_n+\Lambda_n)^{\dagger}z\rangle}.
$$

\subsection{Honesty of confidence intervals}

\begin{theorem}
Suppose Assumptions (A1), (A2), (A3), (B1), (B2) hold. Suppose also that $\lambda_n$ satisfies Eq.~(\ref{eq:lambdan-1}).If 
\begin{equation}
\frac{\frac{1}{n}\sum_{i=1}^n |g(Z_1,Z_i)|^3}{\sqrt{n}(\frac{1}{n}\sum_{i=1}^n g(Z_1,Z_i)^2)^{3/2}}\;\;\overset{p}{\to}\;\; 0
\label{eq:anikrr-lyapnov}
\end{equation}
where $g(z,Z_i)=\langle z,(T_n+\Lambda_n)^{\dagger}Z_i\rangle$ and $Z_i=L_{K^{1/2}}X_i$, then it holds that
\begin{equation}
[\hat\sigma_n^\ani(Z_1)]^{-1}(\hat\eta_n^\ani(X_1)-\eta_0(X_1)) \;\;\overset{d}{\to}\;\;\mathcal N(0,1),
\label{eq:anikrr-clt}
\end{equation}
Furthermore, if $\{\varepsilon_i\}_{i=1}^n$ are Normally distributed, then Eq.~(\ref{eq:anikrr-clt}) holds for \emph{uniformly} for all $X_1,\cdots,X_n$ without the need to satisfy Eq.~(\ref{eq:anikrr-lyapnov}).
\label{thm:anikrr-limiting}
\end{theorem}

\begin{corollary}
Impose the same conditions in Theorem \ref{thm:anikrr-limiting}. Then 
$$
\liminf_{n\to\infty}\inf_{\beta_0\in\Theta_2(\mH)}\Pr[\eta_0(X_1)\in\varphi^\ani(X_1)] \;\; = \;\; 1-\alpha.
$$
where $X_1$ is the first input function in the sample $\{(X_i,Y_i)\}_{i=1}^n$ on which estimation $\hat\eta_n^{\ani}$ and confidence intervals $\varphi^\ani$.
Furthermore, if $\{\varepsilon_i\}_{i=1}^n$ are Normally distributed, the above result can be strengthened to
$$
\liminf_{n\to\infty}\min_{1\leq i\leq n}\inf_{\beta_0\in\Theta_2(\mH)}\Pr[\eta_0(X_i)\in\varphi^\ani(X_i)] \;\; = \;\; 1-\alpha.
$$
without the need to satisfy Eq.~(\ref{eq:anikrr-lyapnov}).
\label{cor:anikrr-ci}
\end{corollary}

\begin{remark}
In both statements of Theorem \ref{thm:anikrr-limiting} and Corollary \ref{cor:anikrr-ci}, we use the notation of $X_1$ and $Z_1$ without loss of generality.
Because $X_1,\cdots,X_n$ are i.i.d., $X_1,Z_1$ can be interpreted as an arbitrary input function in the sample $\{X_i,Y_i\}_{i=1}^n$.
\end{remark}

\subsection{Rates of convergence}

\begin{theorem}
Impose the same conditions in Theorem \ref{thm:anikrr-limiting}. Then the following holds almost surely as $n\to\infty$:
$$
\frac{1}{n}\sum_{i=1}^n \big|\varphi^\ani(X_i)\big|^2 \lesssim \sum_{j=1}^n \frac{\sqrt{\hat s_j}}{\sqrt{\hat s_j}+\sqrt{\lambda_n}}.
$$
where $X_1,\cdots,X_n$ are input functions on which $\hat\eta_n^\ani$ and $\varphi^\ani$ are estimated.
Furthermore, if $\hat s_j\lesssim j^{-2r}$ holds, then the right-hand side of the above inequality is upper bounded by $O(n^{-\frac{2r-1}{2r}+\delta})$
for any $\delta>0$.
\label{thm:anikrr-rates}
\end{theorem}

\subsection{Discussion of limitations}

{\color{red} Discuss the limitations and mention truncation is needed.}

\section{Truncated anisotropic kernel regression}\label{sec:anitr}

In this section we consider \emph{truncated} anisotropic kernel regression estimators, by setting $\lambda_{nj}\equiv +\infty$
for all $j> J(n)$ for a certain truncation level $J(n)\in[n]$ that depends on $n$.
This effectively mutes all eigenfunctions $\{\hat\psi_j\}_{j>J(n)}$, and the resulting estimator $\hat\theta_n$ admits the following spectral expression
\begin{equation}
\hat\theta_n^\anitr = \frac{1}{n}\sum_{i=1}^n\sum_{j=1}^{J(n)}\frac{Y_i'\langle Z_i',\hat\psi_j\rangle}{\hat s_j+\lambda_{nj}}\hat\psi_j.
\label{eq:anitr}
\end{equation}
For a compressed operator notation, define finite-rank operators $T_J$ and $\Lambda_J$ as
$$
T_J := \sum_{j=1}^{J(n)}\hat s_j\hat\psi_j\otimes\hat\psi_j,\;\;\;\;\;\; \Lambda_J := \sum_{j=1}^{J(n)}\lambda_{nj}\hat\psi_j\otimes\hat\psi_j.
$$
Then $\hat\theta_n^\anitr$ can be represented as
$$
\hat\theta_n^\anitr = (T_J+\Lambda_J)^{\dagger}\frac{1}{n}\sum_{i=1}^n Y_i'Z_i',
$$
where $C^\dagger$ is the Moore-Penrose pseudo-inverse of a not necessarily invertible linear operator $C$.

Following the previous section, one option for the choice of the anisotropic Ridge penalty parameters is $\lambda_{nj}=\sqrt{\hat s_j\lambda_n}$ with $\lambda_n$ satisfying Eq.~(\ref{eq:lambdan-1}),
and truncate at the level of $J(n)$ satisfying
\begin{equation}
J(n)n^{-\frac{1}{2r}}\to\infty,\;\;\;\;\;\;J(n)n^{-\frac{1}{2r}-\delta}\to 0
\label{eq:J1}
\end{equation}
as $n\to\infty$, for any $\delta>0$.
The estimator is then $\hat\eta_n^\anitr(\cdot)=\hat\gamma_n + \langle\cdot,\hat\beta_n^\anitr\rangle$ where $\hat\beta_n^\anitr = L_{K^{1/2}}\hat\theta_n^\anitr$.
Afterwards, for a test function $x\in\supp(P_X)$, a confidence interval for $\eta_0(x)$ is constructed as
$$
\varphi^\anitr(x) = [\hat\eta_n^\anitr(x)\;\pm\; \tau_{\alpha}\hat\sigma_n^\anitr(L_{K^{1/2}}x) + \omega_n\|\hat\Pi_J^\perp (L_{K^{1/2}}x)\|_2],
$$
where $\tau_{\alpha}$ is the $(1-\alpha/2)$ quantile of the standard Normal distribution, $\hat\Pi_J^\perp z=(I-\sum_{j=1}^{J(n)}\hat\psi_j\otimes\hat\psi_j)z$, and $\hat\sigma_n^\anitr$ is defined as
$$
\hat\sigma_n^\anitr(z) = ??
$$

\subsection{Additional assumptions and more adaptivity}

\begin{enumerate}
\item[(C1)] There exists a numerical constant $B_1<+\infty$ such that $\sum_{j=1}^{\infty}|\langle \theta_0, \psi_j\rangle| \leq B_1$, where $\{\psi_j\}_{j=1}^{\infty}$ are eigenfunctions of the population covariance operator $T$,
and $\theta_0=L_{K^{-1/2}}\beta_0$ for some $\beta_0\in\mH$.
\end{enumerate}
We use $\Theta_1(\mH)\subseteq\Theta_2(\mH)$ to denote all unknown linear functionals in $\Theta_2(\mH)$ that also satisfies Assumption (D1), for a pre-determined constant upper bound $B_1$.
This assumption is also implicitly imposed, albeit in a different form, in the works of \cite{liu2025scalable} to construct pointwise confidence intervals for classical kernel estimators.

{\color{red}Maybe some remarks on Fourier Basis, or alignment arising from Gaussian processes.}

Let $\mu_n>0$ be a sequence of penalty parameters that decay slightly slower than $\lambda_n$, satisfying
\begin{equation}
\mu_n n^{2r/(2r+1)}\to 0,\;\;\;\;\;\;\mu_n n^{2r/(2r+1)-\delta}\to\infty
\label{eq:mun}
\end{equation}
for all $\delta>0$.
We consider $\hat\theta_n^\adptr$ as solution to Eq.~(\ref{eq:anitr}), with $\Lambda_n=\sum_{j=1}^{J(n)}\lambda_{nj}(z)\hat\psi_j\otimes\hat\psi_j$ being defined as
$$
\lambda_{nj}(z) = \left\{\begin{array}{ll}\hat s_j\sqrt{\mu_n}/|\langle z,\hat\psi_j\rangle|,&\text{if } \langle z,\hat\psi_j\rangle^2\geq\hat s_j;\\ \mu_n,& \text{otherwise};\end{array}\right.
$$
and truncation occurs at $J(n)=\lceil n^{\frac{4r}{4r^2-1}+\delta}\rceil$ for a small $\delta>0$.
Note that in this estimator, the choice of $\Lambda_n$ and the penalization parameters $\lambda_{nj}$ are \emph{adaptive} to the testing function $z=L_{K^{1/2}}x$,
which is different from the previous estimators that do not depend on $z$.

For a testing function $z$ and point estimate $\hat\eta_n^\adptr(x)=\hat\gamma_n+\langle L_{K^{1/2}}z,\hat\theta_n^\adptr\rangle$, a confidence interval of $\eta_0(x)$ is constructed as
$$
\varphi^\adptr(x) = [\hat\eta_n^\adptr(x)\;\pm\;\tau_{\alpha}\hat\sigma_n^\adptr(L_{K^{1/2}}x)+??],
$$
where $\tau_{\alpha}$ is the $(1-\alpha/2)$ quantile of the standard Normal distribution, and $\hat\sigma_n^\adptr$ is defined as
$$
\hat\sigma_n^\adptr(z) = ??
$$


\subsection{Spectrum stability of the sample covariance operator}

In this section, we establish rigorously the connections between the (leading) spectrum of $T$ and $T_n$,
with cutoff thresholds set at levels that can be as deep as $o(\sqrt{n})$. Such discussion is important for the analysis of truncated anisotropic kernel regression estimators, for two primary reasons:
\begin{enumerate}
\item To establish asymptotic Normality for the widest class of testing functions $x\in\supp(P_X)$ possible, we must use the central limit theorem for independently but not identically distributed random variables
and verify classical conditions such as the Lyapnov condition or the Lindeberg condition. This involves understanding the distributions of ``influence functions'' in the variance term, which must be connected to
their population variants in order to avoid statistical correlation between $Z_i$ and $T_n$;
\item The condition (C1) introduced in the previous section, unlike the condition (A2) that $\int_I \theta_0(t)^2\ud t\leq 1$, is \emph{not} basis invariant. This means that on a different orthonormal basis,
such as the eigenfunctions $\{\hat\psi_j\}_{j=1}^n$ of the sample covariance operator $T_n$, the sum $\sum_{j\geq 1}|\langle Z_j,\hat\psi_j\rangle|$ may not be bounded by $B$,
which leads to sub-optimal convergence rates if such sums indeed scale polynomially with $n$.
Therefore, we must rigorously compare the spectrum $\{\psi_j\}$ and $\{\hat\psi_j\}$ of the population and sample covariance operators,
to ensure that $\sum_{j\geq 1}|\langle Z_j,\hat\psi_j\rangle|$ remains bounded so that the bias of kernel Ridge regression can be effectively controlled.
\end{enumerate}

To establish such spectrum stability, we need the following assumption which lower bounds spectral gaps between consecutive eigenvalues on the population covariance operator:
\begin{enumerate}
\item[(D1)] There exists a numerical constant $c_g>0$ such that for all $j\geq 1$, $s_{j+1}-s_j\geq c_g j^{-2r-1}$, where $\{s_j\}_{j=1}^{\infty}$ are eigenvalues (listed in descending order) of the population covariance operator $T=\mathbb E_{P_X}[L_{K^{1/2}}X\otimes L_{K^{1/2}}X]$.
\end{enumerate}

Recall that the population covariance operator $T=\mathbb E[Z\otimes Z]$ and the sample covariance operator $T_n=\frac{1}{n}\sum_{i=1}^nZ_i\otimes Z_i$ admit spectral decomposition of
$T=\sum_{j=1}^{\infty}s_j\psi_j\otimes\psi_j$ and $T_n=\sum_{j=1}^n\hat s_j\hat\psi_j\otimes\hat\psi_j$, where $\{\psi_j\}_{j=1}^{\infty}$ is an orthonormal basis of $L_2(I)$ and $\{\hat\psi_j\}_{j=1}^n$ are orthonormal in $L_2(I)$.
\begin{lemma}
Suppose Assumptions (A1), (A3), (B1), (B2), (B3') and (D1) hold, with $\kappa<1/4$.
Suppose $J(n) \asymp n^{q}$ for some $q\in(0,1/2)$. 
Then there exist constants $C_J,L_J<+\infty$ depending only on $q$ and other parameters in assumptions, such that for every $n\geq 1$, 
with probability $1-n^{-1}$ the following holds:
\begin{enumerate}
\item For every $j\leq J(n)$, there exists $j'\leq C_J j\leq C_JJ(n)$ such that $\|\hat\psi_j-\psi_{j'}\|_2\leq L_J j\sqrt{\log n}/\sqrt{n}$ and $|\hat s_j/s_{j'}-1|\leq L_J j\sqrt{\log n}/\sqrt{n}$;
additionally, for every $k\in[C_J J(n)]\backslash\{j'\}$, $|\langle\hat\psi_j,\psi_k\rangle|\leq L_J\Delta_\psi(j',k)$;
\item For every $j'\leq J(n)/C_J$, there exists $j\leq C_Jj'\leq J(n)$ such that $\|\hat\psi_j-\psi_{j'}\|_2\leq L_J j'\sqrt{\log n}/\sqrt{n}$ and $|\hat s_j/s_{j'}-1|\leq L_J j'\sqrt{\log n}/\sqrt{n}$;
additionally, for every $k\in[C_J J(n)]\backslash\{j'\}$, $|\langle\hat\psi_j,\psi_k\rangle|\leq L_J\Delta_\psi(j',k)$.
\end{enumerate}
In the above statements, the function $\Delta_\psi(j,k)$ admits the following expression:
$$
\Delta_\psi(j,k) = \sqrt{\frac{\log n}{n}}\times \left\{\begin{array}{ll} k^r/j^r,& \text{if }k< j/C;\\ j^r/k^r,& \text{if } k> Cj;\\ \sqrt{jk}/|j-k|;& \text{if $j/C \leq k \leq Cj$ and $j\neq k$.}\end{array}\right.
$$
Furthermore, if $q<1/4$ then the both properties can be strengthened to $\|\hat\psi_j-\psi_j\|_2\leq \epsilon/J(n)$ and $|\hat s_j/s_j-1|\leq \epsilon/J(n))$ for all $j\leq J(n)$,
with the understanding that $j'\equiv j$.
\label{lem:stability}
\end{lemma}

{\color{red} Mention why this can't be done using Davis Kahan, citing existing fPCA papers.}

\subsection{Honesty of confidence intervals}

We first consider the estimator and the confidence interval $\hat\eta_n^\anitr$ and $\varphi^\anitr$. The following theorem establishes the limiting distribution of $\hat\eta_n^\anitr$.
\begin{theorem}
Impose all assumptions and conditions in Lemma \ref{lem:stability}. Suppose $\Lambda_n$ is constructed as $\lambda_{nj}=\sqrt{\hat s_j\lambda_n}$ for all $j\leq J(n)$,
where $\lambda_n,J(n)$ satisfy Eqs.~(\ref{eq:lambdan-1},\ref{eq:J1}). Suppose also that Assumption (B2) holds. Then for any $x\in\supp(P_X)$ and $z=L_{K^{1/2}}x$, 
$$
\hat\sigma_n^\anitr(z)^{-1}(\hat\eta_n^{\anitr}(x)-\eta_0(x)) \;\;\overset{d}{\to}\;\;\mathcal N(0, 1).
$$
\label{thm:anitr-limiting}
\end{theorem}

\subsection{Technical lemmas for results in Sec.~\ref{sec:anitr}.}

\begin{lemma}
Let $\{\psi_j\}_{j=1}^{\infty}$ and $\{\hat\psi_j\}_{j=1}^n$ be the eigenfunctions of $T$ and $T_n$, listed in descending order with respect to their eigenvalues.
Let $J=J(n)$ and $\und J=J(n)/C_J$.
Conditioned on the results in Lemma \ref{lem:stability}, it holds that
$$
\|\hat\Pi_{J}^\perp z\|_2^2 \leq \|\Pi_{\und J}^\perp z\|_2^2 + L_J\times \sum_{j=1}^{\und J}\sum_{k\neq j}^J\langle z,\psi_k\rangle^2\Delta_\psi(j,k)^2\;\;\;\;\;\;\forall z\in L_2(I),
$$
where $\hat\Pi_{J}^\perp = I-\sum_{j=1}^{J}\hat\psi_j\otimes\hat\psi_j$ and $\hat\Pi_{\und J}^\perp = I-\sum_{j=1}^{\und J}\psi_j\otimes\psi_j$ are self-adjoint linear projection operators,
and the summation of $k$ is over all $\{1,2,\cdots,J\}$ except $j$. Furthermore, with $X\sim P_X$ and $Z=L_{K^{1/2}}X$, it holds that
$$
\mathbb E[\|\hat\Pi_J^\perp Z\|_2^2] \lesssim ??.
$$
\label{lem:truncation}
\end{lemma}

\begin{proof}{Proof of Lemma \ref{lem:truncation}.}
Fix arbitrary $z\in L_2(I)$ and, without loss of generality, assume $\|z\|_2=1$.
The Pythagorean theorem asserts that
\begin{align}
\|z\|_2^2 = \|\hat\Pi_{J}z\|_2^2+\|\hat\Pi_{J}^\perp z\|_2^2 = \|\Pi_{\und J}z\|_2^2+\|\Pi_{\und J}^\perp z\|_2^2,
\label{eq:proof-truncation-1}
\end{align}
where $\hat\Pi_J=\sum_{j=1}^J\hat\psi_j\otimes\hat\psi_j$, $\Pi_{\und J}=\sum_{j=1}^{\und J}\psi_j\otimes\psi_j$ and $\hat\Pi_J^\perp=I-\hat\Pi_J$, $\Pi_{\und J}^\perp = I-\Pi_{\und J}$.
Invoking Lemma \ref{lem:stability}, for every $j\leq \und J$, there exists $\ell\leq J$ such that $|\langle\hat\psi_\ell,\psi_j\rangle| = 1-o(1)\leq 1$, and
$$
\big|\langle\hat\psi_\ell, \psi_k\rangle\big| \leq L_J\Delta_\psi(j,k),\;\;\;\;\;\;\forall k\in[J]\backslash\{j\}.
$$
Subsequently,
\begin{align}
\|\Pi_{\und J}z\|_2^2 &= \sum_{j=1}^{\und J}\langle z,\psi_j\rangle^2
\end{align}
\end{proof}

\subsection{Proof of Theorem \ref{thm:anitr-limiting}}

\subsection{Rates of convergence}

%Let $\mH\subseteq L_2(I)$ be a sufficiently smooth subset of all square-integrable functions on $L_2(I)$. For example, $\mH$ could be the Sobolev class with smoothness level $\rho>1/2$.
%Let $X_1,\cdots,X_n$ be i.i.d.~supported on $\mH$, and $\varepsilon_1,\cdots,\varepsilon_n\overset{i.i.d.}{\sim}\mathcal N(0,1)$. Is it true that, with probability $1-n^{-1}$,
%$$
%\sup_{\psi\in\mH}\frac{\frac{1}{n}\sum_{i=1}^n \langle X_i,\psi\rangle^2-\mathbb E[\langle X,\psi\rangle^2]}{\sqrt{\mathbb E[\langle X,\psi\rangle^4]}} \lesssim \gamma(\mH)\sqrt{\frac{\log n}{n}}?
%$$
%And
%$$
%\sup_{\psi\in\mH}\frac{\frac{1}{n}\sum_{i=1}^n \varepsilon_i\langle X_i,\psi\rangle}{\sqrt{\mathbb E[\langle X,\psi\rangle^2]}}\lesssim \gamma(\mH)\sqrt{\frac{\log n}{n}}?
%$$
%Is there a reference for both inequalities? What is the form of $\gamma(\mH)$ and how does it depend on $\rho$? If $1/\sqrt{n}$ is not possible, what is the right rate?
 
\subsection{Minimax optimality of $\varphi^{\adptr}$}


\section{Proofs}

For brevity we assume $\mathbb E[x]=0$, $\gamma_0=0$ and $X_i'=X_i$ throughout all proofs, which is conventional in the analysis of linear functional regression in the literature \citep{cai2012minimax,yuan2010reproducing}.
For all \emph{lower bound} proofs, we also assume $L_{K^{1/2}}=L_{K^{-1/2}}$ is the identity operator to simplify the proofs.
For notational brevity, we many times drop the identity operator symbol $I$, such as writing $T_n+\lambda$ for $T_n+\lambda I$, whenever there is no ambiguity.



\subsection{Some technical lemmas for results in Sec.~\ref{sec:isotropic}}

We first establish a few technical lemmas which are important to our proof of later results.
\begin{lemma}
Suppose Assumptions (A1), (A3) and (B3) hold, and $\lambda_n$ satisfies Eq.~(\ref{eq:lambdan-1}). Then for any $\delta\in(0,1/2]$ there exists a constant $n_0\in\mathbb N$ depending only on $\delta$ and other parameters in assumptions,
 such that for all $n\geq n_0$, with probability at least $1-n^{-1}$,
$$
(1-\delta)(T+\lambda_n)\preceq T_n+\lambda_n \preceq (1+\delta)(T+\lambda_n)
$$
where $T_n=\frac{1}{n}\sum_{i=1}^n Z_i\otimes Z_i$, $T=\mathbb E[Z\otimes Z]$, $Z_i=L_{K^{1/2}}X_i$, $Z=L_{K^{1/2}}Z$, and for two self-adjoint operators $A,B$, $A\preceq B$ if $\langle x,Ax\rangle\leq \langle x, Bx\rangle$
for any square integrable on function $x:I\to\mathbb R$.
\label{lem:empirical-population-cov}
\end{lemma}

We remark that similar versions of Lemma \ref{lem:empirical-population-cov} have appeared, in different forms, in classical analysis of kernel ridge regression estimators and RKHS approaches for linear functional regression.
See e.g.~the upper bound on $\mathcal S_2(\lambda,z)$ in \citep{caponnetto2007optimal}, the second inequality in the proof of Theorem 5 in \citep{smale2007learning}, or Lemma 2 in \citep{cai2012minimax}.
Nevertheless, in our applications the choice of $\lambda_n$ is significantly smaller than conventional choices (for inference purposes), and when $\lambda_n = o(n^{-1})$ existing arguments cannot be directly applied,
without additional conditions such as Assumption (B3) for some $\kappa<1/2$. Therefore, we give the complete proof of Lemma \ref{lem:empirical-population-cov} below for our purposes.

\begin{proof}[Proof of Lemma \ref{lem:empirical-population-cov}.]
We use Hoeffding's inequality for bounded self-adjoint operators, as developed and stated in \citep{minsker2017some}.
For every $i\in[n]$, define
$$
W_i := (T+\lambda_n)^{-1/2}Z_i = (T+\lambda_n)^{-1/2}L_{K^{1/2}}Z_i.
$$
This gives rise to
\begin{align}
\frac{1}{n}\sum_{i=1}^n W_i\otimes W_i &= (T+\lambda_n )^{-1/2}T_n (T+\lambda_n )^{-1/2};\label{eq:epcov-0}\\
\mathbb E[W_i\otimes W_i] &= (T+\lambda_n)^{-1/2} T (T+\lambda_n )^{-1/2}, \;\;\;\;\;\;\forall i\in[n]\label{eq:epcov-0half}.
\end{align}

By Assumption (B3), it holds almost surely that 
\begin{align}
\|W_i\otimes W_i\|_\op&=\left\|\sum_{j,k=1}^{\infty}\frac{\langle Z_i,\psi_j\rangle\langle Z_i,\psi_k\rangle}{\sqrt{s_j+\lambda_n}\sqrt{s_k+\lambda_n}}\psi_j\otimes\psi_k\right\|_\op\nonumber\\
&\leq \max_{j,k\geq 1}s_j^{-1/2}s_k^{-1/2}|\langle Z_i,\psi_j\rangle|\langle Z_i,\psi_k\rangle|\leq C_\nu n^{2\kappa}.
\label{eq:proof-epcov-1}
\end{align}

Let 
$$
A_n := \sum_{i=1}^n\mathbb E[(W_i\otimes W_i)^2] =n\times\sum_{j,k=1}^{\infty}\frac{\mathbb E[\langle Z,Z\rangle\langle Z,\psi_j\rangle\langle Z,\psi_k\rangle]}{(s_j+\lambda_n)(s_k+\lambda_n)}\psi_j\otimes \psi_k.
$$
It is easy to verify that $\|A_n\|_\op \geq n\mathbb E[\langle Z,\psi_1\rangle^4]/(s_1+\lambda_n)=ns_1^2/(s_1+\lambda_n)\geq c_sn/2$ for sufficiently large $n$. 
On the other hand, Because $\langle Z,Z\rangle\leq 1$ almost surely thanks to Assumption (A3), it holds that
\begin{equation}
\tr(A_n) \leq n\times \sum_{j=1}^{\infty} \frac{s_j}{(s_j+\lambda_n)^2} \lesssim n^{(2r+1)/r},
\label{eq:proof-epcov-2}
\end{equation}
where the last inequality holds because $s_j\asymp j^{-2r}$ thanks to Assumption (A1), and $\lambda_n\gtrsim n^{-2}$.
Invoking \citep[Theorem 3.1]{minsker2017some}, it holds with probability $1-n^{-1}$ that
\begin{align*}
\left\|\frac{1}{n}\sum_{i=1}^n W_i\otimes W_i-\mathbb E[W_i\otimes W_i]\right\|_\op &\lesssim n^{-\frac{1-2\kappa}{2}}\log n = o(1). 
\end{align*}
Subsequently, incorporating Eqs.~(\ref{eq:epcov-0},\ref{eq:epcov-0half}), the above inequality implies
$$
\left\|(T+\lambda_n)^{-1/2}(T-T_n)(T+\lambda_n)^{-1/2}\right\|_\op = o(1).
$$
This proves Lemma \ref{lem:empirical-population-cov}, by selecting $n_0$ sufficiently large such that the right-hand side of the above inequality is upper bounded by $\delta$ for any $\delta>0$.
\end{proof}

For the purpose of the next lemma, for any $z=L_{K^{1/2}}x$ for some $x\in\supp(P_X)$ and $Z_i=L_{K^{1/2}}X_i$, define
$$
g(z,Z_i) := \langle z, (T_n+\lambda_n)^{-1}Z_i\rangle,\;\;\;\;\;\;\bar g(z,Z_i) := \langle z, (T+\lambda_n)^{-1}Z_i\rangle.
$$
The following lemma derives probabilistic convergence between $g$, $\bar g$, and their expectations.
\begin{lemma}
Under the same set of conditions in Lemma \ref{lem:empirical-population-cov}, for every $x\in\supp(P_X)$ and $z=L_{K^{1/2}}x$ it holds that
\begin{align*}
\frac{1}{n}\sum_{i=1}^n [g(z,Z_i)-\bar g(z,Z_i)]^2\;\;&\overset{p}{\to}\;0;\\
\frac{1}{n}\sum_{i=1}^n \bar g(z,Z_i)^2-\mathbb E[\bar g(z,Z)^2]\;\;&\overset{p}{\to}\; 0.
\end{align*}
\label{lem:gdiff}
\end{lemma}
\begin{proof}[Proof of Lemma \ref{lem:gdiff}.]
Throughout this proof, we write $w := (T+\lambda)^{-1/2} z$ and $W_i := (T_n+\lambda)^{-1/2}X_i$ for notational simplicity, where $\lambda=\lambda_n$.
Define $\Delta_i := g(z,Z_i)-\bar g(z,Z_i)$. Then
\begin{align*}
\Delta_i &= \langle z,((T_n+\lambda )^{-1}-(T+\lambda )^{-1})Z_i\rangle\nonumber\\
&= \langle z, (T+\lambda )^{-1}((T+\lambda )(T_n+\lambda )^{-1}-I)Z_i\rangle\nonumber\\
&= \langle z, (T+\lambda )^{-1}(T_n-T)(T_n+\lambda )^{-1}Z_i\rangle.
\end{align*}
Subsequently,
\begin{align}
\frac{1}{n}\sum_{i=1}^n \Delta_i^2 &= \langle z, (T+\lambda)^{-1}(T_n-T)(T_n+\lambda)^{-1}T_n(T_n+\lambda)^{-1}(T_n-T)(T+\lambda)^{-1}z\rangle\nonumber\\
&= \langle w, (T+\lambda)^{-1/2}(T_n-T)(T_n+\lambda)^{-1}T_n(T_n+\lambda)^{-1}(T_n-T)(T+\lambda)^{-1/2}w\rangle\nonumber\\
&\leq \langle w, (T+\lambda)^{-1/2}(T_n-T)(T_n+\lambda)^{-1}(T_n-T)(T+\lambda)^{-1/2}w\rangle\nonumber\\
&\leq \|(T+\lambda)^{-1/2}(T_n-T)(T_n+\lambda)^{-1/2}\|_\op\|w\|_2^2\;\; \overset{p}{\to}\;\; 0,
\end{align}
where the last convergence in probability holds by invoking Lemma \ref{lem:empirical-population-cov} and noting that $w=(T+\lambda)^{-1/2}z$ is deterministic. This proves the first property.

For the second identity, because $\{g(z,Z_i)^2\}_{i=1}^n$ are i.i.d.~random variables,
the Chebyshev's inequality asserts that for any $t>0$, 
\begin{align*}
\Pr\left[\left|\frac{1}{n}\sum_{i=1}^n \bar g(z,Z_i)^2 - \mathbb E[g(z,Z)^2]\right| > t\mathbb E[g(z,Z)^2]\right] \leq \frac{\mathbb E[g(z,Z)^4]}{nt^2(\mathbb E[g(z,Z)^2])^2} \leq \frac{C_p}{nt^2},
\end{align*}
where the last inequality holds thanks to Assumption (B2). This implies
\begin{equation*}
\frac{1}{n}\sum_{i=1}^n\bar g(z,Z_i)^2 \;\;\overset{p}{\to}\;\; \mathbb E[g(z,Z)^2],
\end{equation*} 
which completes the proof.
\end{proof}

\begin{lemma}
Assume the same set of conditions in Lemma \ref{lem:empirical-population-cov} holds. 
Let $\varepsilon_i = Y_i - \langle X_i,\beta_0\rangle$. 
Then for any $x\in\supp(P_X)$ and $z=L_{K^{1/2}}x$, as $n\to\infty$, 
$$
\frac{1}{\tilde\sigma_n(z)\sqrt{n}}\sum_{i=1}^n \varepsilon_i g(z,Z_i)\;\;\overset{d}{\to}\;\; \mathcal N(0,1),
$$
where $Z_i=L_{K^{1/2}}X_i$, $g(z,Z)=\langle z,(T_n+\lambda_n I)^{-1}Z\rangle$ and $\tilde\sigma_n(z)^2 = n^{-1}\sum_{i=1}^n \sigma_\varepsilon^2g(z,Z_i)^2$.
\label{lem:isoclt}
\end{lemma}

\begin{proof}[Proof of Lemma \ref{lem:isoclt}.]
Define $\bar g(z,Z) := \langle z,(T+\lambda_n I)^{-1}Z\rangle$.
Because $\{\varepsilon_i\bar g(z,Z_i)\}_{i=1}^n$ are i.i.d.~centered random variables with finite variance, the standard central-limit theorem yields that
\begin{equation}
\frac{1}{\bar\sigma(z)\sqrt{n}}\sum_{i=1}^n \varepsilon_i\bar g(z,Z_i) \;\;\overset{d}{\to}\;\;\mathcal N(0,1),
\label{eq:proof-isoclt-1}
\end{equation}
where
$$
\bar\sigma(z) =\sigma_\varepsilon^2\mathbb E[g(z,Z)^2]= \sigma_\varepsilon^2\mathbb E[\langle z, (T+\lambda_n I)^{-1}Z\rangle^2]
$$

Using Cauchy-Schwarz inequality, it holds almost surely that
\begin{align*}
\left|\frac{1}{n}\sum_{i=1}^n g(z,Z_i)^2-\bar g(z,Z_i)^2\right| &\leq \left|\frac{2}{n}\sum_{i=1}^n [g(z,Z_i)-\bar g(z,Z_i)]\bar g(z,Z_i)\right|\\
&\leq 2\sqrt{\frac{1}{n}\sum_{i=1}^n[g(z,Z_i)-\bar g(z,Z_i)]^2}\sqrt{\frac{1}{n}\sum_{i=1}\bar g(z,Z_i)^2}.
\end{align*}
Lemma \ref{lem:gdiff} then asserts that the first term on the right-hand side of the above inequality converges in probability to zero, and the second term converges in probability to $\sqrt{\mathbb E[g(z,Z_i)^2]}$.
Therefore,
\begin{equation}
\tilde\sigma_n(z)^2 \;\;\overset{p}{\to}\;\; \bar\sigma(z)^2.
\label{eq:proof-isoclt-2}
\end{equation}

Next, using Chebyshev's inequality, for any $\delta>0$, conditioned on $\{Z_i\}$ it holds with probability $1-\delta$ that
$$
\left|\frac{1}{\sqrt{n}}\sum_{i=1}^n\varepsilon_i[g(z,Z_i)-\bar g(z,Z_i)]\right| \leq \frac{\sigma_\varepsilon^2}{\delta^2}\sqrt{\frac{1}{n}\sum_{i=1}^n[g(z,Z_i)-\bar g(z,Z_i)]^2} 
$$
Because $\frac{1}{n}\sum_{i=1}^n[g(z,Z_i)-\bar g(z,Z_i)]^2\overset{p}{\to}0$, for any $\epsilon,u>0$ there exists $N\in\mathbb N$ such that for any $n\geq N$, $\frac{1}{n}\sum_{i=1}^n[g(z,Z_i)-\bar g(z,Z_i)]^2\leq\epsilon$
with probability $1-u$.
Now for any $\epsilon',\delta'>0$, pick $\delta=\delta'/2$, $u=\delta'/2$ and $\epsilon=(\epsilon')^2\delta^4/\sigma_\varepsilon^4>0$, there exists $N\in\mathbb N$ such that for any $n\geq N$, 
$$
\Pr\left[\left|\frac{1}{\sqrt{n}}\sum_{i=1}^n\varepsilon_i[g(z,Z_i)-\bar g(z,Z_i)]\right|>\epsilon'\right]\leq\delta',
$$
where the probability is on both $\{\varepsilon_i\}$ and $\{Z_i\}$, and the above inequality holds by the union bound.
This implies that
\begin{equation}
\frac{1}{\bar\sigma(z)\sqrt{n}}\sum_{i=1}^n\varepsilon_i[g(z,Z_i)-\bar g(z,Z_i)] \;\;\overset{p}{\to}\;\; 0.
\label{eq:proof-isoclt-3}
\end{equation}
Combining Eqs.~(\ref{eq:proof-isoclt-1},\ref{eq:proof-isoclt-2},\ref{eq:proof-isoclt-3}), the proof of Lemma \ref{lem:isoclt} is complete with the Slutsky's theorem.
\end{proof}


\subsection{Proof of Theorem \ref{thm:isokrr-limiting}}

To simplify notations, in the proofs of Theorems \ref{thm:isokrr-limiting} and \ref{thm:isokrr-rates} we shall drop all $\iso$ superscripts as this is the only set of estimators, confidence intervals and their related
quantities to be analyzed in the proofs.
With $Y_i=\langle X_i,\beta_0\rangle+\varepsilon_i = \langle Z_i,\theta_0\rangle+\varepsilon_i$, it holds that
$$
\hat\theta_n = (T_n+\lambda_n)^{-1}T_n\theta_0 + \frac{1}{n}\sum_{i=1}^n\varepsilon_i (T_n+\lambda_n)^{-1}Z_i.
$$
Subsequently,
\begin{align}
\hat\theta_n-\theta_0 &= \underbrace{-\lambda_n (T_n+\lambda_n)^{-1}\theta_0}_{=: b_n} + \underbrace{\frac{1}{n}\sum_{i=1}^n\varepsilon_i(T_n+\lambda_n)^{-1}Z_i}_{=:v_n}.
\end{align}

For any test function $x\in\supp(P_X)$, let $z=L_{K^{1/2}}x=\sum_{j=1}^{n}\hat\nu_j\hat\psi_j$, where $\hat\nu_j=\langle z,\hat\psi_j\rangle$.
Let also $\hat\vartheta_j = \langle\theta_0,\hat\psi_j\rangle$.
Slightly abusing notations by denoting $b_n(z) := \langle b_n,z\rangle = \int_I b_n(t)z(t)\ud t$, we have that
\begin{align}
\big|b_n(z)\big| &= \lambda_n\left|\left\langle(T_n+\lambda_n)^{-1}\sum_{j=1}^n\hat\vartheta_j\hat\psi_j, \sum_{j=1}^n\hat\nu_j\hat\psi_j\right\rangle\right|
= \lambda_n \left|\sum_{j=1}^n \frac{\hat\nu_j\hat\vartheta_j}{\hat s_j+\lambda_n}\right|\label{eq:proof-simple-1}\\
&\leq \lambda_n \sqrt{\sum_{j=1}^n\hat\vartheta_j^2}\sqrt{\sum_{j=1}^n\frac{\hat\nu_j^2}{(\hat s_j+\lambda_n)^2}} \label{eq:proof-simple-2}\\
&\leq \sqrt{\lambda_n}\sqrt{\sum_{j=1}^n\hat\vartheta_j^2}\sqrt{\sum_{j=1}^n\frac{\hat\nu_j^2}{\hat s_j+\lambda_n}}\label{eq:proof-simple-2half}\\
&= \sqrt{\lambda_n}\cdot \sqrt{\langle z,(T_n+\lambda_n)^{-1}z\rangle} = \sqrt{\lambda_nn}\hat\sigma_n(z).\label{eq:proof-simple-3}
\end{align}
In the above derivations, Eq.~(\ref{eq:proof-simple-1}) holds because $\{\hat\psi_j\}_{j=1}^n$ is orthonormal with respect to $L_2(I)$ and $T_n=\sum_{j=1}^n\hat s_j\hat\psi_j\otimes\hat\psi_j$;
Eq.~(\ref{eq:proof-simple-2}) holds by the Cauchy-Schwarz inequality;
Eq.~(\ref{eq:proof-simple-2half}) holds by noting that $\lambda_n/(\sqrt{\hat s_j+\lambda_n})\leq\sqrt{\lambda_n}$ for all $j$, and furthermore
$$
\langle z, (T_n+\lambda_n)^{-1}z\rangle = \left\langle\sum_{j=1}^n\hat\nu_j\hat\psi_j, \sum_{j=1}^n\frac{\hat\nu_j}{\hat s_j+\lambda_n}\hat\psi_j\right\rangle = \sum_{j=1}^n\frac{\hat\nu_j^2}{\hat s_j+\lambda_n}.
$$
With $\lambda_n n\to 0$ as $n\to\infty$, Eq.~(\ref{eq:proof-simple-3}) implies that
$$
b_n(z)/\hat\sigma_n(z)\;\;\to\;\;0\;\;\;\;\;\;a.s.
$$

We next analyze the term involving $v_n$ on an arbitrary test function $x\in\supp(P_X)$ with $z=L_{K^{1/2}}x=\sum_{j=1}^n\hat\nu_j\hat\psi_j$, where $\hat\nu_j=\langle z,\hat\psi_j\rangle$.
Let $v_n(z)=\langle v_n,z\rangle = \int_I v_n(t)z(t)\ud t$. We have that
$$
v_n(z) = \frac{1}{n}\sum_{i=1}^n \varepsilon_i \langle z, (T_n+\lambda_n)^{-1}Z_i\rangle.
$$
Note that $v_n(z)$ is \emph{pivotal}, meaning that it does not depend on the unknown $\beta_0\in\Theta_2(\mH)$ or $\theta_0=L_{K^{1/2}}\beta_0$.
Note also that $\{\varepsilon_i\}_{i=1}^n$ are conditionally centered when $\{Z_i\}_{i=1}^n$ are fixed. Therefore, if $\{\varepsilon_i\}_{i=1}^n$ are Normally distributed,
the distribution of $v_n(z)$ is an exact Normal distribution with mean zero and variance equal to
$$
\frac{1}{n^2}\sum_{i=1}^n \langle z,(T_n+\lambda_n)^{-1}Z_i\rangle^2 = \frac{1}{n}\langle z, (T_n+\lambda_n)^{-1}T_n(T_n+\lambda_n)^{-1}z\rangle = \frac{\tilde\sigma_n(z)^2}{n}
$$
because $T_n=\frac{1}{n}\sum_{i=1}^n Z_i\otimes Z_i$. If $\{\varepsilon_i\}_{i=1}^n$ are not Normally distributed, Lemma \ref{lem:isoclt} asserts that under imposed conditions,
$$
\frac{1}{\tilde\sigma_n(z)\sqrt{n}}\sum_{i=1}^n\varepsilon_i\langle z,(T_n+\lambda_n)^{-1}Z_i\rangle\;\;\overset{d}{\to}\;\; \mathcal N(0,1),
$$
which is to be proved.

\subsection{Proof of Theorem \ref{thm:isokrr-rates}}

For any $x\in\supp(P_X)$ and $z=L_{K^{1/2}}x$, the definition of $\varphi(x)$ yields
$$
n\big|\varphi(x)\big|^2 = \sigma_\varepsilon^2\langle z, (T_n+\lambda_n)^{-1}z\rangle^2.
$$
Because $|\varphi(x)|$ is pivotal, we may safely drop the $\sup_{\beta_0\in\Theta_2(\mH)}$ notation and instead calculate
\begin{align*}
n\mathbb E_{P_X}\left[|\varphi(X)|^2\right] &= n\mathbb E_{P_Z}[\langle Z,(T_n+\lambda_n)^{-1}Z\rangle^2] = \sigma_\varepsilon^2\tr((T_n+\lambda_n)^{-1/2}T(T_n+\lambda_n^{-1/2}))\\
&= \sigma_\varepsilon^2(1+o_P(1))\tr((T+\lambda_n)^{-1/2}T(T+\lambda_n)^{-1/2}),
\end{align*}
where the last identity holds by invoking Lemma \ref{lem:empirical-population-cov}. Additionally, invoking Assumption (A1) and the fact that $\lambda_n\gtrsim n^{-1-\delta}$ for any $\delta>0$, it holds that
\begin{align*}
\tr((T+\lambda_n)^{-1/2}T(T+\lambda_n^{-1/2})) &= \sum_{j=1}^{\infty} \frac{s_j}{s_j+\lambda_n} \lesssim \int_1^{\infty}\frac{j^{-2r}}{j^{-2r}+\lambda_n}\ud j\lesssim n^{1/2r+\delta},
\end{align*}
which is to be proved.

\subsection{Proof of Theorem \ref{thm:anikrr-limiting}}

In the proofs of Theorems \ref{thm:anikrr-limiting} we shall drop all $\ani$ superscripts as this is the only set of estimators, confidence intervals and their related quantities to be analyzed in these proofs.
Following similar derivations that lead to Eq.~(\ref{eq:proof-isoclt-1}), and accounting for the null space of $\{\hat\psi_j\}_{j=1}^n$, we have for the anisotropic case that
\begin{equation}
\hat\theta_n-\theta_0 =\underbrace{-(I-\hat\Psi)\theta_0}_{=:t_n}\underbrace{- (T_n+\Lambda_n)^{\dagger}\Lambda_n\theta_0}_{=: b_n} + \underbrace{\frac{1}{n}\sum_{i=1}^n\varepsilon_i(T_n+\Lambda_n)^{\dagger}Z_i}_{=:v_n},
\label{eq:proof-aniclt-1}
\end{equation}
where $\hat\Psi = \sum_{j=1}^n\hat\psi_j\otimes\hat\psi_j$.

Consider $X_1$ being an arbitrary input function in $\{(X_i,Y_i)\}_{i=1}^n$,
and let $Z_1=L_{K^{1/2}}X_1$, $b_n(Z_1)=\langle b_n,Z_1\rangle$, $v_n(Z_1)=\langle v_n,Z_1\rangle$.
Noet that $t_n(Z_1)=0$ because 
Because $v_n(z)$ is either exactly Normally distributed (if $\{\varepsilon_i\}_{i=1}^n$ are Normally distributed), or satisfies Lyapnov conditions of central limit theorem with probabilities converging to one
provided that Eq.~(\ref{eq:anikrr-lyapnov}) is satisfied, we know that $v_n(z)$ is asymptotically Normally distributed whose variance is exactly $\frac{\sigma_\varepsilon^2}{n^2}\sum_{i=1}^n\langle z, (T_n+\Lambda_n)^{-1}Z_i\rangle^2 = \frac{\sigma_\varepsilon^2}{n}\langle z, (T_n+\Lambda_n)^{-1}T_n(T_n+\Lambda_n)^{-1}z\rangle=\hat\sigma_n(z)^2$. This shows
\begin{equation}
\frac{v_n(z)}{\hat\sigma_n(z)}\;\;\overset{d}{\to}\;\;\mathcal N(0,1).
\label{eq:proof-aniclt-2}
\end{equation}

We next analyze the bias term. Using the spectral decomposition of both $T_n$ and $\Lambda_n$, $b_n(z)$ can be expressed as
$$
b_n(z) = \left\langle(T_n+\Lambda_n)^{-1}\Lambda_n\sum_{j=1}^n\hat\vartheta_j\hat\psi_j, \sum_{j=1}^n\hat\nu_j\hat\psi_j\right\rangle = \sum_{j=1}^n\frac{\lambda_{nj}\hat\nu_j\hat\vartheta_j}{\hat s_j+\lambda_{nj}},
$$
where $\hat\vartheta_j=\langle\theta_0,\hat\psi_j\rangle$ and $\hat\nu_j = \langle z,\hat\psi_j\rangle$.
Using Cauchy-Schwarz inequality and the fact that $\sum_{j=1}^n\hat\vartheta_j^2\leq 1$ because $\|\theta_0\|_2^2 = \int_I\theta_0(t)^2\ud t\leq 1$, it holds that
$$
b_n(z)^2 \leq \sum_{j=1}^n\frac{\lambda_{nj}^2\hat\nu_j^2}{(\hat s_j+\lambda_{nj})^2}.
$$
Using the choice that $\lambda_{nj}=\sqrt{\hat s_j\lambda_n}$, the above inequality simplifies to
\begin{align}
b_n(z)^2 \leq \lambda_n\sum_{j=1}^n\frac{\hat\nu_j^2}{(\sqrt{\hat s_j}+\sqrt{\lambda_n})^2}.
\label{eq:proof-aniclt-3}
\end{align}
On the other hand, expanding $\hat\sigma_n^2(z)$ in the spectrum of $\{\hat\psi_j\}_{j=1}^n$ we obtain
\begin{align}
\hat\sigma_n^2(z) &= \frac{\sigma_\varepsilon^2}{n}\langle z, (T_n+\Lambda_n)^{-1}T_n(T_n+\Lambda_n)^{-1}z\rangle = \frac{\sigma_\varepsilon^2}{n}\sum_{j=1}^n\frac{\hat\nu_j^2\hat s_j}{(\hat s_j+\lambda_{nj})^2}\nonumber\\
&=  \frac{\sigma_\varepsilon^2}{n}\sum_{j=1}^n\frac{\hat\nu_j^2}{(\sqrt{\hat s_j}+\sqrt{\lambda_n})^2}.
\label{eq:proof-aniclt-4}
\end{align}
Comparing Eqs.~(\ref{eq:proof-aniclt-3},\ref{eq:proof-aniclt-4}) and noting that $\lambda_n n\to\infty$ as $n\to\infty$, we obtain
\begin{equation}
\frac{|b_n(z)|}{\hat\sigma_n(z)} \leq \sigma_\varepsilon\sqrt{\lambda_n n} \to 0\;\;\;\;\;a.s.
\label{eq:proof-aniclt-5}
\end{equation}
Eqs.~(\ref{eq:proof-aniclt-1},\ref{eq:proof-aniclt-5}) together yield Theorem \ref{thm:anikrr-limiting}, by Slutsky's theorem.

\subsection{Proof of Theorem \ref{thm:anikrr-rates}}

Let $Z_i=L_{K^{1/2}}X_i$.
It holds that
\begin{align*}
\frac{1}{n}\sum_{i=1}^n\big|\varphi(X_i)\big|^2 &= \frac{4\tau_\alpha^2}{n^2}\sum_{i=1}^n\hat\sigma_n(z_i)^2 = \frac{4\sigma_\varepsilon^2\tau_\alpha^2}{n^2}\sum_{i=1}^n \langle Z_i, (T_n+\Lambda_n)^{-1}T_n(T_n+\Lambda_n)^{-1}Z_i\rangle\\
&= \frac{4\sigma_\varepsilon^2\tau_\alpha^2}{n}\tr(((T_n+\Lambda_n)^{-1}T_n)^2) = \frac{4\sigma_\varepsilon^2\tau_\alpha^2}{n}\sum_{j=1}^n\frac{\hat s_j^2}{(\hat s_j+\lambda_{nj})^2}\\
&\leq \frac{4\sigma_\varepsilon^2\tau_\alpha^2}{n}\sum_{j=1}^n\frac{\hat s_j}{\hat s_j+\lambda_{nj}} = \frac{4\sigma_\varepsilon^2\tau_\alpha^2}{n}\sum_{j=1}^n\frac{\sqrt{\hat s_j}}{\sqrt{\hat s_j}+\sqrt{\lambda_{n}}}\\
&\lesssim \frac{1}{n}\int_1^{\infty} \frac{j^{-r}}{j^{-r}+\sqrt{\lambda_n}}\ud j \lesssim n^{-\frac{2r-1}{2r}+\delta}
\end{align*}
for any $\delta>0$, where the last inequality holds by the condition imposed on $\hat\{s_j\}$ and the asymptotic scaling that $\lambda_n\gtrsim n^{-1-\delta}$ for any $\delta>0$.


\subsection{Proof of Lemma \ref{lem:stability}}

Express both $T$ and $T_n$ on the first $n$ eigenfunctions $\psi_1,\cdots,\psi_n$. That is, we construct $n\times n$ matrices $A$ and $A_n$, such that 
$A_{ij}=\langle\psi_i, T\psi_j\rangle$ and $[A_n]_{ij}=\langle\psi_i, T_n\psi_j\rangle$.
It is easy to verify that $A$ is a diagonal matrix with diagonal entries $s_1,\cdots,s_n$, while $A_n$ is a positive semi-definite matrix with a more nuanced structure.
More specifically,
$$
A_{ij} = \mathbb E[\langle Z,\psi_i\rangle\langle Z,\psi_j\rangle] = \vct 1\{i=j\}s_i.
$$
On the other hand, $[A_n]_{ij}$ is the sample average of a sum of $n$ i.i.d.~random variables:
$$
[A_n]_{ij} = \frac{1}{n}\sum_{\ell=1}^n \langle Z_\ell,\psi_i\rangle\langle Z_\ell,\psi_j\rangle.
$$

%{\color{blue}
%Consider the submatrix of $A_n$ corresponding to eigenfunctions $J(n)+1,\cdots,n$. The trace formula and Eq.~(\ref{eq:proof-stability-1}) implies that, with probability $1-n^{-1}$,
%\begin{align}
%\sum_{j=J(n)+1}^n\hat s_j &= \sum_{j=J(n)+1}^n [A_n]_{jj} \leq \sum_{j=J(n)+1} s_j + O\left(\frac{s_j\sqrt{\log n}}{\sqrt{n}} + \frac{s_jn^{2\kappa}\log n}{n}\right)\nonumber\\
%&\lesssim \int_{J(n)}^{\infty} j^{-2r}\left(1+\frac{\sqrt{\log n}}{\sqrt{n}}\right)\ud j  \lesssim J(n)^{-2r+1}.
%\end{align}
%This proves the second property of Lemma \ref{lem:stability}.
%}

%{\color{blue}
%Consider the submatrix of $A_n$ corresponding to eigenfunctions $J(n)+1,\cdots,n$. The trace formula yields
%\begin{align}
%\sum_{j=J(n)+1}^n\hat s_j-s_j &= \sum_{j=J(n)+1}^{\infty} [A_n]_{jj}-A_{jj} = \sum_{j=J(n)+1}^n \langle\psi_j, (T_n-T)\psi_j\rangle.
%\label{eq:proof-stability-1-1}
%\end{align}
%For every $i\in[n]$, define $w_i := \sum_{j=J(n)+1}^{\infty}\langle\psi_j, Z_i\rangle^2$ such that $\mathbb E[w_i] = \sum_{j=J(n)+1}^{\infty}s_j$.
%The right-hand side of Eq.~(\ref{eq:proof-stability-1-1}) is precisely $\frac{1}{n}\sum_{i=1}^n w_i-\mathbb E[w_i]$.
%Furthermore, 
%\begin{align*}
%\mathbb E[w_i^2] &=\sum_{j,k>J(n)}\mathbb E[\langle Z_i,\psi_j\rangle^2\langle Z_i,\psi_k\rangle^2]\overset{(a)}{\leq} \left(\sum_{j>J(n)}\sqrt{\mathbb E[\langle Z_i,\psi_j\rangle^4]}\right)^2\\
%&\overset{(b)}{\leq} \left(\sum_{j>J(n)}C_ps_j\right)^2 \overset{(c)}{\lesssim} \left(\int_{J(n)}^{\infty} j^{-2r}\ud j\right)^2 \lesssim J(n)^{2-4r},
%\end{align*}
%where Eq.~(a) holds by the Cauchy-Schwarz inequality, Eq.~(b) holds by Assumption (B2) and Eq.~(c) holds by Assumption (B1).
%Additionally, Assumption (A3) implies that $w_i\leq \|Z\|_2^2\leq 1$ almost surely. Subsequently, using Bernstein's inequality, with probability $1-n^{-1}$ it holds that
%\begin{align}
%\left|\sum_{j=J(n)+1}^n \langle\psi_j, (T_n-T)\psi_j\rangle\right| &= \left|\frac{1}{n}\sum_{i=1}^n w_i-\mathbb E[w_i]\right| \lesssim \frac{J(n)^{1-2r}\sqrt{\log n}}{\sqrt{n}} + \frac{\log n}{n}.
%\label{eq:proof-stability-1-2}
%\end{align}
%Combining Eqs.~(\ref{eq:proof-stability-1-1},\ref{eq:proof-stability-1-2}) and noting that $\sum_{j>J(n)}s_j \lesssim \int_{J(n)}^{\infty} j^{-2r}\ud j \lesssim J(n)^{1-2r}$, we complete the proof of the second property of Lemma \ref{lem:stability}.
%}


Now, let $C=(2C_s/3c_s)^{1/(2r)}>0$ be a constant such that, for every $j\in\mathbb N$, $s_{Cj}\leq C_s (Cj)^{-2r} \leq c_sj^{-2r}/2\leq s_j/2$ thanks to Assumption (B1). 
Let $\bar J=\lceil CJ(n)\rceil$.
Restrict both $A$ and $A_n$ to their $\bar J\times \bar J$ sub-matrices corresponding to the first $J(n)$ eigenfunctions, $\psi_1,\cdots,\psi_{\bar J}$.
That is, both $A$ and $A_n$ are now $\bar J\times\bar J$ positive semi-definite matrices.
By Assumption (B2) and the Cauchy-Schwarz inequality, it holds that 
$$
\mathbb E[\langle Z,\psi_i\rangle^2\langle Z,\psi_j\rangle^2] \leq \sqrt{\mathbb E[\langle Z,\psi_i\rangle^4]}\sqrt{\mathbb E[\langle Z,\psi_j\rangle^4]}
\leq C_p\mathbb E[\langle Z,\psi_i\rangle^2]\mathbb E[\langle Z,\psi_j\rangle^2] = C_ps_is_j.
$$
By Assumption (B3'), it holds that 
$$
|\langle Z,\psi_i\rangle\langle Z,\psi_j\rangle| \leq C_\nu^2\sqrt{s_js_k}n^{2\kappa}
$$ almost surely.
Subsequently, using Bernstein's inequality and the union bound (Bonferroni inequalities), with probability $1-n^{-1}$ it holds uniformly over $i,j\in[\bar J]$ that 
\begin{align}
\big|[A_n]_{ij}-A_{ij}\big| \leq C_1\left(\sqrt{\frac{C_ps_js_k\log n}{n}} + \frac{C_\nu\sqrt{s_js_k}n^{2\kappa}\log n}{n}\right) \lesssim \sqrt{\frac{s_js_k\log n}{n}},
\label{eq:proof-stability-1}
\end{align}
where $C_1\leq 5$ is a numerical constant, and the last inequality holds provided that $\kappa<1/4$.

Let $e_1,e_2,\cdots$ be the standard basis vectors.
We will prove that, for every $j\in[\bar J]$, there exists $w\in\mathbb R^{\bar J}$ satisfying $\|w\|_2 = O(j\sqrt{\log n}/\sqrt{n})=o(1)$, $w_j=0$ and $|w_k|\lesssim \Delta_\psi(j,k)$ for all $k\in[\bar J]\backslash\{j\}$ such that
\begin{equation}
A_n (e_j+w) = s_j(e_j+w),
\label{eq:eigen-an}
\end{equation}
and furthermore the convergence in the $o(1)$ term to zero as $n\to\infty$ is \emph{uniform} over $j\in[\bar J]$.
This shows that the self-adjoint operator $T_n=\sum_{k,\ell=1}^{n}[A_n]_{k\ell}\psi_k\otimes\psi_\ell$ admits a pair of eigenvalue and eigenfunction $(\hat s,\hat\psi)$ that satisfies
$\hat s=s_j(1+O(j\sqrt{\log n}/\sqrt{n}))=s_j(1+o(1))$, $\|\hat\psi-\psi_j\|_2=O(j\sqrt{\log n}/\sqrt{n})=o(1)$ and $|w_k|\lesssim \Delta_\psi(j,k)$ for all $k\neq j$. Then
\begin{enumerate}
\item For every $j\leq J(n)$, with all $j'\in\{1,2,\cdots,C j\}$ these $(\hat s,\hat\psi)$ pairs (a total number of $C j$ such pairs) cover all the first $j$ eigenvalues and eigenfunctions
of $T_n$, which proves the first property of Lemma \ref{lem:stability}.
\item For every $j\leq J(n)$, with $\hat s_\ell\geq (1+o(1))s_{C j}\geq (1+o(1))c_s (C_Jj)^{-2r}\geq (1+o(1))C_s (j/C)^{-2r}$ for all $\ell\leq j$, 
the $\{\hat s_\ell,\hat\psi_\ell\}_{\ell=1}^{j}$ pairs cover all the first $j/C$ eigenvalues and eigenfunctions of $T$, which proves the second property of Lemma \ref{lem:stability}.
\end{enumerate}

Let $H=A_n-A$ and consider $w_j=0$ which is simply a Gauge choice. Because $Ae_j=s_je_j$, Eq.~(\ref{eq:eigen-an}) implies
$$
(A-s_j I+H)w = -He_j.
$$
Let $\bar H\in\mathbb R^{(\bar J-1)\times (\bar J-1)}$ be the submatrix formed by removing the $j$th row and $j$th column of $H$. Let $h\in\mathbb R^{\bar J-1}$ be the $j$th column of $H$, removing $H_{jj}$.
Let $D=\diag(s_1-s_j,\cdots,s_{j-1}-s_j,s_{j+1}-s_j,\cdots,s_{\bar J}-s_j)\in\mathbb R^{(\bar J-1)\times(\bar J-1)}$. Let $\bar w=(w_1,\cdots,w_{j-1},w_{j+1},\cdots,w_n)\in\mathbb R^{\bar J-1}$. Then
$$
\bar w = -(D+\bar H)^{-1}h = -D^{-1/2}(I+\tilde H)^{-1}D^{-1/2}h,
$$
where $\tilde H = D^{-1/2} \bar H D^{-1/2}$. 

For any $i,k\neq j$, Eq.~(\ref{eq:proof-stability-1}) implies
$$
\big|\tilde H_{ik}\big| \lesssim \sqrt{\frac{\log n}{n}} \frac{i^{-r}k^{-r}\sqrt{\log n}}{\sqrt{|(s_i-s_j)(s_k-s_j)|}} =: \sqrt{\frac{\log n}{n}}z_iz_k\;\;\;\;\;\;\text{where}\;\; z_k = \frac{k^{-r}}{\sqrt{|s_k-s_j|}}.
$$
In a similar way, $D^{-1/2}h$ can also be upper bounded component-wise as
\begin{equation}
\big|[D^{-1/2}h]_k \big| \lesssim  \sqrt{\frac{\log n}{n}}\frac{j^{-r}k^{-r}}{\sqrt{|s_k-s_j|}} = \frac{j^{-r}\sqrt{\log n}}{\sqrt{n}}z.
\label{eq:proof-stability-3}
\end{equation}

Discuss the following three cases:
\begin{enumerate}
\item The case of $k\leq j/C$. In this case, by Assumption (B1), $|s_k-s_j|\geq c_s k^{-2r} - C_s(Ck)^{-2r}\geq 0.5c_sk^{-2r}$, and therefore $|z_k|\leq \sqrt{2/c_s}=O(1)$;
\item The case of $k\geq Cj$. In this case, by Assumption (B1), $|s_k-s_j| \geq c_s j^{-2r} - C_s (Cj)^{-2r}\geq 0.5c_s j^{-2r}$, and therefore $|z_k|\leq \sqrt{2/c_s}j^r/k^r= O(j^rk^{-r})$.
\item The case of $j/C\leq k\leq Cj$. In this case, by Assumption (D1), $|s_k-s_j| \geq |k-j| \cdot c_g(Cj)^{-2r-1}$, and therefore $|z_k|\leq \sqrt{C}j^{r+1/2}k^{-r}/\sqrt{c_g|k-j|}\leq C^{r+1}\sqrt{c_gk/|j-k|}= O(\sqrt{k/|j-k|})$.
\end{enumerate}
Because $\tilde H=\sqrt{\frac{\log n}{n}}zz^\top$, summarizing the above cases we obtain 
\begin{align}
\|\tilde H\|_\op&\lesssim \sqrt{\frac{\log n}{n}}\|z\|_2^2 =\sqrt{\frac{\log n}{n}}\sum_{k\neq j}|z_k|^2  \lesssim \frac{1}{\sqrt{n}}\left(j + j^{2r}\int_j^{\infty}t^{-2r}\ud t + Cj\int_1^{Cj} \frac{\ud t}{t}\right)\nonumber\\
& \lesssim \frac{j\ln (Cj+1)\sqrt{\log n}}{\sqrt{n}}\lesssim \frac{j\log^2 n}{\sqrt{n}}.
\end{align}
With $\bar J\asymp n^q$ for $q\in(0,1/2)$ (Because $J(n)\asymp n^q$ and $\bar J=CJ$, where $C$ is a constant), the right-hand side of the above inequality is $o(1)$ and therefore we have $\|\tilde H\|_\op=o(1)$ as $n\to\infty$.
We can then do the following Taylor expansion using $(I+\tilde H)^{-1} = \sum_{\ell=0}^{\infty} (-1)^\ell \tilde H^\ell$, for every $k\neq j$:
\begin{align}
|w_k| &\leq \sum_{\ell=0}^{\infty} \big|e_k^\top D^{-1/2}\tilde H^\ell D^{-1/2}h\big| \overset{(a)}{\leq }\sum_{\ell=0}^{\infty} \frac{j^{-r}\sqrt{\log n}}{\sqrt{n}}\left|e_k^\top D^{-1/2}\left(\sqrt{\frac{\log n}{n}}zz^\top\right)^\ell z\right|\nonumber\\
&\lesssim \sum_{\ell=0}^{\infty}\frac{j^{-r}\sqrt{\log n}}{\sqrt{n}}\left(\frac{j\log^2 n}{\sqrt{n}}\right)^\ell \big|e_k^\top D^{-1/2}z\big| \lesssim \frac{j^{-r}\sqrt{\log n}}{\sqrt{n}}\big|e_k^\top D^{-1/2}z\big|\nonumber\\
& = \frac{j^{-r}\sqrt{\log n}}{\sqrt{|s_k-s_j| n}}\big|z_k\big|.\label{eq:proof-stability-4} 
\end{align}
Here, Eq.~(a) holds by plugging in the expression of $D^{-1/2}h$ in Eq.~(\ref{eq:proof-stability-3}).
Invoking the case discussions earlier. For $k\leq j/C$, the right-hand side of Eq.~(\ref{eq:proof-stability-4}) is upper bounded by
$$
 \frac{j^{-r}\sqrt{\log n}}{\sqrt{|s_k-s_j|n}}\times O(1) = O\left(\frac{j^{-r}k^r\sqrt{\log n}}{\sqrt{n}}\right);
$$
for $k\geq Cj$, the right-hand side of Eq.~(\ref{eq:proof-stability-4}) is upper bounded by 
$$
\frac{j^{-r}\sqrt{\log n}}{\sqrt{|s_k-s_j|n}}\times O(j^rk^{-r}) = O\left(\frac{j^rk^r\sqrt{\log n}}{\sqrt{n}}\right);
$$
for $j/C\leq k\leq Cj$, the right-hand side of Eq.~(\ref{eq:proof-stability-4}) is upper bounded by
\begin{align*}
\frac{j^{-r}\sqrt{\log n}}{\sqrt{|s_k-s_j|n}}\times O\left(\sqrt{\frac{k}{|j-k|}}\right) &= \frac{j^{-r}\sqrt{\log n}}{\sqrt{n}}\times O\left(\frac{j^{r+1/2}}{\sqrt{|k-j|}}\right)\times O\left(\sqrt{\frac{k}{|j-k|}}\right)\\
&= O\left(\frac{\sqrt{jk\log n}}{|j-k|\sqrt{n}}\right).
\end{align*}
Combining all of the above cases we complete the proof of upper bounds on $|\langle\hat\psi_j,\psi_k\rangle|$ in Lemma \ref{lem:stability}.

Furthermore, 
\begin{align}
\|w\|_2 &\leq \frac{j^{-r}\sqrt{\log n}}{\sqrt{n}}\left(\int_1^j \frac{\ud t}{t^{-2r}} + j^{2r}\int_j^{\bar J}\frac{t^{-2r}}{j^{-2r}}\ud t + Cj\cdot \int_1^{Cj}\frac{1}{tj^{-2r-1}}\frac{\ud t}{t}\right)^{1/2}\nonumber\\
&\lesssim \frac{j^{-r}\sqrt{\log n}}{\sqrt{n}}\left(j^{2r+1} + j^{2r+1} + j^{2r+2}\right)^{1/2} \lesssim  \frac{j\sqrt{\log n}}{\sqrt{n}}\leq \frac{\bar J\sqrt{\log n}}{\sqrt{n}}.\label{eq:proof-stability-5}
\end{align}
Because $\bar J\asymp n^q$ for some $q<1/2$, the right-hand side of Eq.~(\ref{eq:proof-stability-5}) is $o(1)$ uniformly for all $j\leq\bar J$. This completes the proof.

Finally, with the additional assumption that $q<1/4$, the right-hand side of Eq.~(\ref{eq:proof-stability-5}) is further upper bounded by $o(1/\bar J)=o(1/J(n))$. 
Additionally, in this case we have for every $j\leq \bar J$ that $\hat s_j-\hat s_{j+1} = s_j-s_{j+1}+o(s_j/\bar J)=s_j-s_{j+1}+o(j^{-2r}/\bar J) = (1+o(1))(s_j-s_{j+1})$ because $s_j-s_{j+1}\gtrsim j^{-2r-1}$.
This shows that there is a one-to-one match between $(s_j,\psi_j)$ and $(\hat s_j,\hat\psi_j)$ for all $j\leq\bar J$.

\bibliographystyle{apalike}
\bibliography{refs}

\end{document}